\chapter{ارتباط با کدِ بیرونی}%
\label{ch:ffi}

دنیای نرم‌افزار پر از کتابخانه‌های آماده است که در زبانِ \lr{C} نوشته شده‌اند:
از کارکردهای ریاضیِ سیستم تا پایگاه‌های داده و کتابخانه‌های گرافیکی. سلام به شما
اجازه می‌دهد مستقیماً این کتابخانه‌ها را صدا بزنید. به این پل، \textbf{رابطِ
کارکردِ بیرونی} (FFI) می‌گویند.

\section{فراخوانیِ یک کارکردِ بیرونی}

برای صدا زدنِ یک کارکردِ \lr{C}، نخست با یک بلوکِ \code{بیرونی} (معادلِ
\code{extern}) به سلام \emph{معرفی}اش می‌کنیم یعنی می‌گوییم نامش چیست و چه
ورودی و خروجی دارد. این مثالِ واقعی چند کارکردِ ریاضیِ کتابخانهٔ استانداردِ
\lr{C} را صدا می‌زند:

\begin{salamen}
extern:
    func printf(format : str, ...) : int
    func sqrt(x : f64) : f64
    func pow(base : f64, exp : f64) : f64
    func hypot(x : f64, y : f64) : f64
end
const PI : f64 = 3.14159265358979
func main:
    printf("sqrt(2)    = %.6f\n", sqrt(2.0))
    printf("2^10       = %.0f\n", pow(2.0, 10.0))
    printf("hypot(3,4) = %.1f\n", hypot(3.0, 4.0))
end
\end{salamen}

بلوکِ \code{extern:} می‌گوید: «این کارکردها در جایی دیگر در کدِ \lr{C} ---
تعریف شده‌اند؛ من فقط امضایشان را اعلام می‌کنم.» سپس می‌توانید آن‌ها را درست
مانندِ کارکردهای سلام صدا بزنید: \code{sqrt(2.0)}، \code{pow(2.0, 10.0)}.

\begin{note}
به \code{...} در امضای \code{printf} دقت کنید. این نشانه می‌گوید که کارکرد
\emph{تعدادِ متغیری} از آرگومان می‌پذیرد. \code{printf} کارکردِ معروفِ چاپ در
\lr{C} است که یک «رشتهٔ قالب» می‌گیرد و سپس هر تعداد مقدار که آن قالب بخواهد.
نشانه‌هایی مانندِ \lr{\texttt{\%.6f}} (عددِ اعشاری با ۶ رقم) و
\lr{\texttt{\%d}} (عددِ صحیح) جای مقدارها را در رشته مشخص می‌کنند.
\end{note}

\section{همان مفهوم به فارسی}

واژهٔ \code{بیرونی} نیز مانندِ همهٔ واژگانِ سلام معادلِ فارسیِ
\code{extern} است. این مثالِ واقعی نشان می‌دهد که بلوکِ بیرونی را می‌توان در یک
برنامهٔ فارسی هم نوشت:

\begin{salamfa}
بیرونی:
    کارکرد sqlite3_libversion() : str
پایان

کارکرد آغازین:
    v : str = sqlite3_libversion()
    اگر v != "":
        چاپ "کتابخانهٔ sqlite متصل شد"
    وگرنه:
        چاپ "اتصال برقرار نشد"
    پایان
پایان
\end{salamfa}

\section{پیوند به یک کتابخانه}

برخی کارکردها (مانندِ \code{printf} و \code{sqrt}) در کتابخانه‌های پایه‌ای‌اند که
به‌طورِ خودکار در دسترس‌اند. اما برای کتابخانه‌های دیگر (مانندِ \code{sqlite})
باید به سلام بگویید که به کدام کتابخانه \emph{پیوند} بزند. این کار با
\code{پیوند پویا} (dynamic link) انجام می‌شود، و معمولاً همراهِ همان دستورِ
شرطیِ \code{اگر} برای انتخابِ نامِ درست در هر سیستم‌عامل، فقط این‌بار در سطحِ
بالای فایل و در زمانِ کامپایل ارزیابی می‌شود:

\begin{salamfa}
اگر SALAM_OS_WINDOWS:
    پیوند پویا "winsqlite3"
وگرنه:
    پیوند پویا "sqlite3"
پایان
بیرونی:
    کارکرد sqlite3_libversion() : str
پایان

کارکرد آغازین:
    چاپ sqlite3_libversion()
پایان
\end{salamfa}

اینجا روی ویندوز به \code{winsqlite3} و روی سیستم‌های دیگر به \code{sqlite3}
پیوند می‌زنیم نمونه‌ای زیبا از همکاریِ FFI و شرطِ زمانِ کامپایل که در فصلِ پیش
دیدیم.

\section{کاربردِ عملی: خواندنِ متغیرِ محیطی از C}

این مثالِ واقعی، کارکردِ \code{getenv} از \lr{C} را برای خواندنِ نامِ کاربر صدا
می‌زند:

\begin{salamen}
extern:
    func printf(format : str, ...) : int
    func getenv(name : str) : str
    func strlen(s : str) : u64
end

func main:
    user : str = getenv("USERNAME")
    printf("Your username is %d bytes.\n", strlen(user))
end
\end{salamen}

اینجا سه کارکردِ \lr{C} را با هم به کار گرفتیم: \code{getenv} (خواندنِ متغیرِ
محیطی)، \code{strlen} (طولِ یک رشتهٔ \lr{C}) و \code{printf} (چاپ). این نشان
می‌دهد که چگونه سلام می‌تواند مستقیم با توانمندی‌های سیستم‌عامل از راهِ \lr{C}
کار کند.

\begin{warn}
کار با FFI قدرتمند اما حساس است. وقتی به کدِ \lr{C} می‌روید، از حفاظ‌های ایمنیِ
سلام بیرون می‌آیید: باید خودتان مطمئن شوید که ریخت‌ها و امضاها دقیقاً درست‌اند.
یک امضای نادرست (مثلاً \code{i32} به‌جای \code{i64}) می‌تواند برنامه را به‌طورِ
ناگهانی از کار بیندازد. همیشه مستنداتِ کتابخانهٔ \lr{C} را با دقت بخوانید.
\end{warn}

\begin{deep}[نگاهی به درون: FFI چگونه کار می‌کند]
وقتی سلام برنامهٔ شما را به \lr{C} ترجمه می‌کند، کارکردهای معمولیِ سلام را با یک
نامِ ویژه «مزیّن» می‌کند تا با هم تداخل نکنند. اما کارکردهای \code{extern} را
\emph{بدونِ} این تزئین و با نامِ اصلیِ \lr{C} صدا می‌زند دقیقاً همان‌طور که
کامپایلرِ \lr{C} انتظار دارد. به همین دلیل است که سلام می‌تواند بی‌واسطه با هر
کتابخانهٔ \lr{C} کار کند: هر دو در نهایت به همان «ABIِ زبانِ C» می‌رسند، که
زبانِ مشترکِ تقریباً همهٔ نرم‌افزارهای سیستمی است.
\end{deep}

\section{خلاصهٔ فصل}

\begin{itemize}
  \item FFI پلی است برای صدا زدنِ کارکردهای \lr{C} از درونِ سلام.
  \item با بلوکِ \code{بیرونی:} (یا \code{extern:}) کارکردهای بیرونی را
        اعلام می‌کنیم.
  \item \code{...} در امضا یعنی کارکرد آرگومان‌های متغیر می‌گیرد (مثلِ
        \code{printf}).
  \item با \code{پیوند پویا "..."} به یک کتابخانه پیوند می‌زنیم.
  \item پیش‌پردازنده به ما کمک می‌کند نامِ درستِ کتابخانه را در هر سیستم‌عامل
        برگزینیم.
  \item FFI شما را از حفاظ‌های ایمنی بیرون می‌برد؛ با دقت کار کنید.
\end{itemize}

\section{تمرین‌ها}

\exercise با FFI، کارکردِ \code{sqrt} از \lr{C} را صدا بزنید و جذرِ چند عدد را
چاپ کنید.

\exercise کارکردِ \code{getenv} را صدا بزنید و مقدارِ متغیرِ محیطیِ \code{PATH}
را بگیرید (می‌توانید فقط طولش را با \code{strlen} چاپ کنید).

\exercise با \code{printf} و رشته‌های قالبِ گوناگون
(\lr{\texttt{\%d}}، \lr{\texttt{\%.2f}}، \lr{\texttt{\%s}})، چند مقدار را با
قالب‌بندیِ زیبا چاپ کنید.

\exercise توضیح دهید چرا امضای نادرستِ یک کارکردِ بیرونی خطرناک است.

\exercise یک بلوکِ \code{بیرونی} بنویسید که چند کارکردِ ریاضیِ \lr{C}
(\code{sin}، \code{cos}، \code{pow}) را اعلام کند و آن‌ها را به کار بگیرد.
